% ====================================================================
% Section 2 (v2): Related work.
% Positions the contribution as parsimony, not SOTA competition.
% ====================================================================
\section{Related Work}
\label{sec:related}

\paragraph{Engineered load balancing.}
Modern multipath load balancers spread traffic to avoid hot spots.
ECMP~\cite{ecmp} splits flows across equal-cost shortest paths by hashing, with
no congestion awareness. CONGA~\cite{conga} adds that awareness: it enumerates a
set of candidate paths and selects among them using in-network congestion
feedback, rerouting at flowlet granularity. DRILL~\cite{drill} pushes the
decision to the switch and the microsecond, making per-packet local choices
based on queue occupancy with minimal state. These systems define the operating
points we measure against. Our aim is not to outperform them---on the
topologies where they are strong, they remain strong---but to ask how much of
their behaviour is recoverable from a far simpler, physically-motivated rule.

\paragraph{Physics- and nature-inspired routing.}
Using physical or biological analogies for routing has a long history:
ant-colony optimisation, electrical-network and resistance analogies,
potential- and field-based forwarding, and physarum-inspired path formation.
These approaches typically introduce a mechanism and demonstrate that it works.
Our contribution differs in emphasis in two ways. First, we run the analogy in
reverse as an ablation: rather than showing that added physics helps, we add a
great deal of physics and then measure how much can be removed, arriving at a
two-term residue. Second, we frame the outcome as a statistical equivalence
against strong engineered baselines, rather than as a demonstration that the
physical method is novel or superior.

\paragraph{Equivalence testing.}
Most networking evaluations test for a difference (one method beats another).
We instead test for \emph{equivalence}, using TOST (two one-sided
tests), which asks whether two methods are indistinguishable within a stated
margin. This is the statistically appropriate tool when the claim is ``as good
as'', not ``better than'', and it guards against the familiar error of reading a
non-significant difference as evidence of equivalence. To our knowledge,
equivalence testing is uncommon in the routing-evaluation literature, and its
use here is deliberate: the entire thesis is one of equivalence with parsimony.

\paragraph{Applicability conditions.}
A recurring weakness of method papers is silence about when the method fails. We
attach to our result an explicit, topology-only predictor (the tube/sp ratio,
Section~\ref{sec:applicability}) that says in advance where the equivalence
holds and where engineered routing retains an edge. This places the
contribution closer to the spirit of complex-systems and network-science work,
where structural predictors of dynamical behaviour are themselves the object of
study, than to a conventional systems-performance paper.


\medskip
Conceptually adjacent, though router-free, is the literature on dynamic
averaging and load-balancing processes on graphs, e.g.\ Alistarh et
al.~\cite{alistarh2022dynamic}, where iterative local exchanges converge
to equilibria fixed by graph structure. The attractor reported in this
paper is the routing-level analogue of that structural determinism,
observed empirically across heterogeneous mechanisms rather than proved
for a single dynamics.
